\section{Envelopes: the common mechanism}\label{sec:envelopes}

Most of our upper-bound-capable methods, with and without a prover, are
instances of one estimator (the exceptions are the correlation-inequality bounds
of \cref{sec:janson,sec:independence}, the product-of-ratios estimators of
\cref{sec:cftp}, and two variants of it: the counter-guided sampler of
\cref{thm:sis}, and the tilted Gibbs estimator of \cref{sec:local}, whose mass and
proposal are themselves certified statistically). It turns a \emph{deductive}
fact (``$C(x)=1$ implies $\varphi(x)\ge1$'') and an \emph{exact} computation
($\E_D\varphi$) into an importance sampler whose cost depends only on how loose
the deduction is.

\subsection{Envelope circuits}

\begin{definition}[envelope]\label{def:envelope}
An \emph{envelope circuit} $\Phi$ is an arithmetic circuit whose leaves are
literal indicators $[x_i=b]$, nonnegative rational constants, and
(\cref{sec:boxes}) \emph{boxes} $[B(x)=1]$ for circuits $B$; whose gates are
$\Sum$ and $\Prod$; and whose $\Prod$ gates are \emph{decomposable}: their children
have pairwise disjoint supports, where the support $\sv(v)$ of a node is the set of
variables occurring at the leaves of its cone ($\sv(B)$ for a box). We also require
that every $\Prod$-child with empty support be a constant leaf (fold constant
subcircuits first). $\Phi$ computes $\varphi:\bits^n\to\R_{\ge0}$. It
\emph{dominates} $C$ if $\varphi(x)\ge1$ whenever $C(x)=1$. Its \emph{mass} is
$Z=\E_D\varphi$; if $p>0$, its \emph{looseness on $C$} is $K=Z/p$, and domination
gives $Z\ge p$, i.e.\ $K\ge1$.
\end{definition}

A \emph{derivation} of $\Phi$ is a proof tree of its tree unfolding: keep,
top-down, one child of every $\Sum$ occurrence and all children of every $\Prod$
occurrence. By decomposability a derivation visits each node of nonempty support
at most once. With literal leaves and unit constants, $\varphi(x)$ is the number of
derivations all of whose leaves are satisfied by $x$; with general constants it is
the corresponding weighted count. A DNF $\bigvee_j T_j$ with consistent,
duplicate-free terms gives the envelope $\sum_j\prod_{\ell\in T_j}[\ell]$, and a DNNF
gives the envelope obtained by reading $\vee$ as $\Sum$ and $\wedge$ as $\Prod$.

\begin{lemma}[exact mass, exact sampling]\label{lem:envmass}
Let $\Phi$ have literal and constant leaves only, $D$ be a product distribution
with rational marginals, and $Z>0$. Then $Z=\E_D\varphi$ is computed by one
bottom-up pass ($\Sum\mapsto$ sum of children's masses, $\Prod\mapsto$ product)
with $O(|\Phi|)$ operations on rationals of $\poly(|\Phi|,b)$ bits ($b$ the total
bit length of the marginals and of the constants of $\Phi$), and the following
top-down procedure samples exactly from $Q_\varphi(x)=D(x)\varphi(x)/Z$: at a
$\Sum$ gate choose a child with probability proportional to its mass; at a
$\Prod$ gate visit all children; at a literal leaf set the literal; finally draw
every unset variable from $D$.
\end{lemma}

\begin{proof}
By decomposability and independence, $\E_D\prod_c\varphi_c=\prod_c\E_D\varphi_c$;
linearity handles $\Sum$. The procedure selects a derivation $\mathcal T$ with
probability $w(\mathcal T)/Z$, where $w(\mathcal T)$ is the product of the
literals' probabilities and the constants in $\mathcal T$ (the choice probabilities
telescope), and then draws $x\sim D$ conditioned on all literals of $\mathcal T$,
whose density is $D(x)\one[x\models\mathcal T]\,c(\mathcal T)/w(\mathcal T)$ with
$c(\mathcal T)$ the product of its constants. Summing over derivations,
$\Pr[x]=\sum_{\mathcal T}c(\mathcal T)D(x)\one[x\models\mathcal T]/Z=D(x)\varphi(x)/Z$.
(A literal is set at most once per derivation: two literals on the same variable
in one derivation would lie below different children of a decomposable $\Prod$.)
Bit lengths: the mass of a node is a sum of at most $2^{|\Phi|}$ products of at most
$|\Phi|$ constants and marginals.
\end{proof}

In floating point, masses must be kept in the log domain and the final interval
widened by the accumulated rounding error. Our implementation does exactly this:
along every derivation the computed weight and the sampler's choice probabilities
telescope, so the relative error of $\hat Z\cdot\E_{\hat Q}W$ is at most
$e^{\pm\eta}$, where $\eta$ is the largest sum of local rounding errors on a
derivation (a max-sum dynamic program over $\Phi$); the discreteness of the
categorical draws at $\Sum$ gates adds an absolute $O((n+|\Phi|)2^{-53})$ to $\E W$
(the remaining inputs are drawn from $D$ exactly, by comparing lazily drawn binary
digits of a uniform with those of $q_i$, so the rejection fills of
\cref{sec:boxes} are exact and the box estimates rest on exact samples); and the
output interval is widened by both and rounded outwards.

\subsection{The envelope estimator}

For $N$ i.i.d.\ samples in $[0,1]$ with empirical mean $\hat\mu$, the
\emph{Hoeffding--KL interval} at level $\delta$ is
$\{m\in[0,1]:N\,\KL(\hat\mu\|m)\le\ln(2/\delta)\}$ with $\KL$ the Bernoulli relative
entropy; it contains the true mean with probability at least $1-\delta$
\cite[Thm.~1]{hoeffding1963}. (In floating point our implementation excludes a value
only when its computed divergence exceeds the bound by more than its rounding error,
and allows for the rounding of $\hat\mu$, so the computed interval contains the exact
one.)

\begin{theorem}[envelope estimator]\label{thm:envelope}
Let $\varphi:\bits^n\to\R_{\ge0}$ dominate $C$, with exact mass $Z$ and an exact
sampler for $Q_\varphi=D\varphi/Z$ (for instance an envelope circuit with
\cref{lem:envmass}). If $Z=0$, output $[0,0]$ (valid under domination). Otherwise,
for $X\sim Q_\varphi$ let $W=C(X)/\varphi(X)$. Then
\begin{enumerate}[label=(\roman*)]
\item $W\in[0,1]$ and $\E W=p/Z=1/K$;
\item from $N$ independent copies, the Hoeffding--KL interval $[\ell,u]$ for
$\E W$ at level $\delta$ gives an interval $[Z\ell,\min(Zu,1)]$ containing $p$ with
probability at least $1-\delta$, \emph{for every circuit $C$} as long as $\varphi$
dominates it;
\item if $\varphi$ does not dominate $C$: for integer-valued $\varphi$ (envelope
circuits with literal leaves and integer constants) the same procedure gives a
valid interval for $\Pr_D[C(x)=1\wedge\varphi(x)>0]\le p$; for general $\varphi$, with
$W$ clipped at $1$, $Z\ell$ is still a valid lower bound on $p$;
\item for $0<\eps\le1$, run sequentially from
$N_0=\lceil c\log(2/\delta_1)/\eps^2\rceil$ samples, doubling $N$ and using level
$\delta_j=6\delta/(\pi^2j^2)$ at stage $j$: all intervals are valid with probability
$1-\delta$, and if $p>0$ the procedure reaches $u\le(1+\eps)\ell$ after
$N=O\big(K(\log(1/\delta)+\log\log K)/\eps^2\big)$ samples with probability $1-\delta$
(so both hold with probability $1-2\delta$). With a sample cap it always halts,
returning its current, still valid, interval. Each sample costs $O(|C|+|\Phi|)$
arithmetic operations with the sampler of \cref{lem:envmass}.
\end{enumerate}
\end{theorem}

\begin{proof}
(i) If $Q_\varphi(x)>0$ then $\varphi(x)>0$; if moreover $C(x)=1$ then
$\varphi(x)\ge1$, so $W\le1$. $\E W=\sum_xD(x)\varphi(x)C(x)/(Z\varphi(x))=p/Z$.
(ii) The KL interval is valid for i.i.d.\ $[0,1]$ variables, and $p\le1$.
(iii) For integer-valued $\varphi$, $\varphi(x)>0$ implies $\varphi(x)\ge1$, so
$W\in[0,1]$ and $\E W=\Pr[C=1,\varphi>0]/Z$. For real-valued $\varphi$, clipping
only lowers the mean, so $\E\min(W,1)\le\Pr[C=1,\varphi>0]/Z\le p/Z$.
(iv) Let $\mu=1/K$. By the multiplicative Chernoff bound, once
$N\ge c\log(1/\delta_j)/(\mu\eps^2)$ the empirical mean $\hat\mu$ is within a factor
$1\pm\eps/8$ of $\mu$ with probability $1-\delta_j$. Since
$\KL(a\|b)\ge(a-b)^2/(2\max(a,b))$, every $m$ in the KL interval with $m>\hat\mu$
satisfies $(m-\hat\mu)^2\le am$ with $a=2\log(2/\delta_j)/N$, hence
$m-\hat\mu\le a+\sqrt{a\hat\mu}$, and similarly below; for a suitable constant $c$
both are at most $(\eps/4)\hat\mu$, so $u/\ell\le(1+\eps/4)/(1-\eps/4)\le1+\eps$. With
the stated $N_0$, the first stage at which this holds has $j=O(\log K)$, whence the
$\log\log K$ term; the union bound over stages makes all intervals simultaneously
valid.
\end{proof}

So the cost of an envelope is governed by one number. Monte Carlo is the
trivial envelope $\varphi\equiv1$, with $K=1/p$; exact computation is $\varphi=C$
(when $C$ is itself a d-DNNF), with $K=1$. In between:

\begin{lemma}[what $K$ measures]\label{lem:K}
For $p>0$, $K=\E_D[\varphi\mid C=1]+\E_D[\varphi\,\one_{C=0}]/p$. For an envelope
with literal leaves and unit constants whose positivity set $\{\varphi>0\}$ equals
$C^{-1}(1)$ (e.g.\ a DNNF for $C$), $K$ is the expected number of satisfied
derivations of $x\sim D\mid C=1$, its \emph{expected ambiguity}. For a DNF with
$m$ terms, $K\le m$.
\end{lemma}

The first term is overcounting (union bounds); the second is mass that the
deduction failed to exclude. $K$ also bounds the relative variance
$\rho=\frac{Z}{p}\E_D[1/\varphi\mid C=1]-1\le K-1$ of $ZW$. The KL procedure of
\cref{thm:envelope}(iv) ignores the variance and pays $\Theta(K)$ samples whatever
$\rho$ is; a variance-adaptive interval (empirical Bernstein, or DKLR's algorithm
$\mathcal{AA}$) pays $O\big((\rho/\eps^2+K/\eps)\log(1/\delta)\big)$, so a low dispersion
of the derivation count helps only in the $\eps^{-2}$ term. The $K/\eps$ term cannot
be avoided by a sound procedure that uses the weights as a black box, and for tilting
envelopes $K$ can exceed $\rho$ exponentially (\cref{prop:certgap}). A d-DNNF has
$\varphi=\one_C$, $K=1$ and $\rho=0$. For $E=\bigwedge_{i\le m}(x_i\vee y_i)$ under the uniform
distribution, used as its own envelope, $K=(4/3)^m$, and the relative variance is at
least $(10/9)^m-1$ under every reweighting of its $\vee$-edges (with equality for unit
weights); both are exponential although $D(E)=(3/4)^m$ is known exactly.
\Cref{thm:sis} removes this overcounting altogether: its looseness is $D(E)/p$ rather
than $\E_D[\varphi]/p$.

In the spirit of surprise accounting \cite{hilton2024surprise}: of the $\log_2(1/p)$
bits of surprise in the event $C(x)=1$, an envelope with $Z\le1$ explains
$\log_2(1/Z)$ bits by deduction and exact computation (once domination is certified,
syntactically or by a proof), and leaves $\log_2K$ bits to sampling, at a cost of
about $K$ samples (up to $\eps^{-2}\log(1/\delta)$ factors). If $Z\ge1$ the envelope is
worse than Monte Carlo and is replaced by $\varphi\equiv1$. Monte Carlo explains
nothing; exact computation explains everything, with $|\Phi|$ possibly exponential
in $|C|$. Surprise accounting also charges the explanation's own description length:
zero for envelopes computed from $C$, and $|\pi|$ for prover-supplied ones
(\cref{thm:surprise}).

\subsection{Dense leaves}\label{sec:boxes}

Products of independent events that are individually not rare are the most
common source of rarity, and they need not have small envelopes: an
$\wedge$ of $k$ independent random sub-circuits $B_j$, each true with
probability $1/2$, has $p=2^{-k}$ but its literal envelope may lose a constant
factor per block. Boxes repair this: a box is a sub-function used as its own
envelope, its mass estimated by Monte Carlo and its conditional distribution
sampled by rejection. Estimation errors in the masses bias the proposal, but
the bias on each derivation is a product of per-box ratios that we can bound
exactly.

\begin{theorem}[envelopes with dense leaves]\label{thm:boxes}
Let $\Phi$ dominate $C$, with literal leaves and box leaves $B_1,\dots,B_r$
($\Prod$ gates decomposable with respect to supports), where
$m_j=D(B_j)\ge1/R$ and every derivation contains at most $L$ box leaves; let
$0<\eps\le1$. Estimate each $m_j$ by Monte Carlo to an interval $[\ell_j,u_j]$ at
level $\delta/(2r)$ with point value $\hat m_j$; compute masses $\hat Z$ and sample
$\hat Q$ as in \cref{lem:envmass} with $\hat m_j$ in place of $m_j$, filling each box
by rejection sampling from $D$. For $\rho(\mathcal T)=\prod_{B_j\in\mathcal
T}\hat m_j/m_j$ let $\rho_{\min},\rho_{\max}$ be its extreme values over derivations
when each $m_j$ ranges over $[\ell_j,u_j]$ (computed by a min-/max-product pass over
$\Phi$). Then, with $[\ell,u]$ a valid KL interval for $\E_{\hat Q}W$, $W=C/\varphi$, at
total level $\delta/2$ (fixed or sequential),
\[
p\in\big[\hat Z\ell/\rho_{\max},\;\hat Zu/\rho_{\min}\big]
\]
with probability at least $1-\delta$. If each box interval satisfies
$u_j\le(1+\eps/(8L))\ell_j$ and importance sampling is run to ratio
$1+\eps/4$, the output is tight with probability $1-2\delta$, after
$O(K(\log(1/\delta)+\log\log K)/\eps^2)$ samples of $\hat Q$, each costing
$O(|C|+|\Phi|+\sum_j|B_j|+LR\max_j|B_j|)$ in expectation; estimating the boxes costs
$O(rL^2R\log(r/\delta)/\eps^2)$ evaluations.
\end{theorem}

\begin{proof}
The sampler picks derivation $\mathcal T$ with probability
$\hat w(\mathcal T)/\hat Z$ where $\hat w$ uses the $\hat m_j$, and then draws
$x$ exactly from $D$ conditioned on $\mathcal T$'s leaves (rejection sampling
is exact; decomposability makes the leaves' conditionings independent), whose
density is $D(x)\one[x\models\mathcal T]/w(\mathcal T)$ with the true masses.
Hence $\hat Q(x)=\frac{D(x)}{\hat Z}\sum_{\mathcal T\ni x}\rho(\mathcal T)$, and
\[
\hat Z\,\E_{\hat Q}W=\sum_xD(x)C(x)\cdot\frac{1}{\varphi(x)}\sum_{\mathcal
T\ni x}\rho(\mathcal T)\;\in\;[\rho_{\min},\rho_{\max}]\cdot p,
\]
because $\varphi(x)$ is the number of derivations satisfied by $x$, so the
middle factor is an average of $\rho(\mathcal T)$. On the event (probability
$\ge1-\delta/2$) that all $m_j\in[\ell_j,u_j]$, $\rho_{\min}$ and $\rho_{\max}$
bound this average; conditionally on the box estimates the KL interval fails
with probability at most $\delta/2$. For tightness: $\rho_{\max}/\rho_{\min}\le
(1+\eps/(8L))^{2L}\le e^{\eps/4}$ and $(1+\eps/4)e^{\eps/4}\le1+\eps$ for $\eps\le1$.
Since $\hat Z=\sum_{\mathcal T}\rho(\mathcal T)w(\mathcal T)\le\rho_{\max}Z$, the mean
of $W$ under $\hat Q$ is at least $\rho_{\min}p/(\rho_{\max}Z)\ge e^{-\eps/4}/K$, so
\cref{thm:envelope}(iv) (which only needs i.i.d.\ $[0,1]$ samples of that mean)
gives the sample count. Computing $\varphi(X)$ evaluates every box, hence the
$\sum_j|B_j|$ term.
\end{proof}

With boxes, a product of $k$ independent dense blocks has $K=1$ and costs
$O(k^3R\log(k/\delta)/\eps^2)$ evaluations, against $2^{\Theta(k)}$ for Monte
Carlo; a DNF whose terms are products of dense blocks has an envelope (a $\Sum$ of
$\Prod$s of boxes) whose positivity set is $C^{-1}(1)$ and with $\varphi\le m$, so
$K\le m$ by \cref{lem:K}.

\subsection{Soft envelopes: Chernoff bounds as certificates}\label{sec:tilting}

Envelopes need not be integer-valued. If
$C(x)=[S(x)\ge t]$ with $S=\sum_{j=1}^k g_j(x_{B_j})$ a sum of
$\{0,1\}$-valued sub-circuits on disjoint blocks (since $S$ is an integer and the
comparison is exact, we take $t$ to be an integer, replacing it by $\lceil t\rceil$),
then for $\theta\ge0$
\[
\varphi_\theta(x)=e^{-\theta t}\prod_{j=1}^k\big([g_j=0]+e^{\theta}[g_j=1]\big)=e^{\theta(S(x)-t)}\ \ge\ \one[S(x)\ge t],\quad
Z_\theta=e^{-\theta t}\prod_{j=1}^k\big(1-p_j+p_je^{\theta}\big),
\]
where $p_j=D(g_j=1)$. $Z_\theta$ is the Chernoff bound, a deductively valid
upper bound, and $Q_{\varphi_\theta}$ is the exponentially tilted product
distribution. To sample it block by block we assume that the $p_j$ are known
exactly and $D$ conditioned on $g_j=b$ can be sampled, e.g.\ blocks small enough
to enumerate (as in our implementation), or boxes with $\min(p_j,1-p_j)\ge1/R$, in
which case \cref{thm:boxes} applies with $L=k$ (with weighted averages of $\rho$).
This is the classical efficient importance sampler for large deviations of sums
\cite{siegmund1976,bucklew2004}; here is a clean uniform bound on its looseness.

\begin{proposition}[tilting is polynomially tight]\label{prop:tilting}
Let the weights be $1$, $\E_D S\le t<k_+$, where $k_+=|\{j:p_j>0\}|$ and $t$ is an
integer. There is $\theta\ge0$ with $\E_{Q_\theta}S=t$, and for it
$K=Z_\theta/p\le\tfrac43\big(2\sqrt{k}+1\big)$. The same bound holds for any
$\theta\ge0$ with $|\E_{Q_\theta}S-t|<1/(k+2)$, so $\theta$ may be found by bisection
and $e^\theta$ taken rational.
\end{proposition}

\begin{proof}
$\theta\mapsto\E_{Q_\theta}S=\sum_jp_je^\theta/(1-p_j+p_je^\theta)$ is continuous and
nondecreasing, at most $t$ at $\theta=0$ and tends to $k_+>t$, so $\theta$ exists.
Under $Q_\theta$ the $g_j$ are independent Bernoulli$(p_j^\theta)$, so $S$ is
Poisson-binomial with mean $\mu_\theta$. By Darroch's rule \cite{darroch1964} the
mode is $t$ when $\mu_\theta=t$, and also when $|\mu_\theta-t|<1/(k+2)$ (the rule
makes $k'$ the unique mode for $k'\le\mu<k'+1/(k'+2)$ and for
$k'-1/(n-k'+2)<\mu\le k'$). By Chebyshev,
$Q_\theta(|S-\mu_\theta|\le2\sigma)\ge3/4$ with $\sigma^2=\mathrm{Var}_{Q_\theta}S\le
k/4$, and this window contains at most $4\sigma+1$ integers, so, $S$ being unimodal
(log-concave), $Q_\theta(S=t)\ge\frac{3}{4(4\sigma+1)}$. Finally
$p=Z_\theta\,\E_{Q_\theta}[e^{-\theta(S-t)}\one_{S\ge t}]\ge Z_\theta\,Q_\theta(S=t)$.
\end{proof}

Our implementation, which also allows nonnegative integer block weights $w_j$, finds
$\theta$ by bisection on $\E_{Q_\theta}S$ to within $1/(2(k+2))$ of $t$, subject to
$\theta\max_jw_j\le699$ (so that every $e^{\theta w_j}$ is a finite double), and compares
$S$ with $\lceil t\rceil$ in integer arithmetic, so $\varphi_\theta\ge1$ wherever $C=1$.

Unit weights matter. For integer weights under a product distribution the best
tilt can be loose by a factor $p^{-\Theta(1)}$: with
$S=(k+1)y+\sum_{i\le k}x_i$, $q_y=\eta$, $q_x=1/2$ and $t=k+1$, so $p=\eta$, the
optimal Chernoff envelope has $K=2^{25.7}$ at $k=20$, $p=2^{-40}$. When the $p_j$ are
known exactly, the unit-weight case is also computable exactly by dynamic
programming over $S$ (pseudo-polynomial in general weights), and arbitrary weights
have a deterministic FPTAS under product distributions via read-once branching programs
\cite{gkm2010} (see also \cite{gopalan2011fptas}). The point of
\cref{prop:tilting} is that the Chernoff bound, a one-line deduction from the syntax
``$C$ is a threshold of independent blocks'', explains all but $O(\log k)$ bits of the
surprise, and sampling pays for the rest. When the blocks overlap, the tilted proposal is
a Gibbs measure; \cref{sec:local} certifies its sampler and normaliser, and
\cref{rem:local-blocks} records that the $O(\sqrt k)$ bound is tight and that $K\to1$
at the top threshold.

\subsection{Certified presumption of independence}\label{sec:janson}

The \emph{presumption of independence} \cite{christiano2022presumption}
estimates $\Pr[\bigwedge_iA_i]$ by $\prod_i\Pr[A_i]$. For monotone events it can
be certified: correlation inequalities bound the error in terms of the
probability carried by dependent pairs, which can be computed or estimated
without ever conditioning on the rare event.

\begin{theorem}[Harris--Janson sandwich]\label{thm:janson}
Let $C=\bigwedge_{i=1}^m\mathrm{cl}_i$ be a CNF in which every variable occurs
with only one sign, $B_i=[\mathrm{cl}_i=0]$ (a subcube), $D$ a product
distribution, $b_i=D(B_i)$, $M=\prod_i(1-b_i)$, $\eps_0=\max_ib_i<1$, and for an
ordering of the clauses $d_i=\sum_{j<i,\,\sv(\mathrm{cl}_j)\cap\sv(\mathrm{cl}_i)\ne\emptyset}D(B_i\wedge B_j)$,
$\Delta=\sum_id_i$ (the sum over unordered dependent pairs, as in
\cref{thm:janson-general}). Then
\[
M\;\le\;p\;\le\;\prod_i(1-b_i+d_i)\;\le\;M\exp\Big(\sum_i\frac{d_i}{1-b_i}\Big)\;\le\;M\exp\Big(\frac{\Delta}{1-\eps_0}\Big),
\]
and all quantities are computable with $O(\sum_i|\mathrm{cl}_i|^2\cdot d)$ arithmetic
operations, $d$ the maximum number of clauses sharing a variable. (If some
$b_i=1$ then $p=M=0$.)
\end{theorem}

\begin{proof}
Flip the variables that occur positively; then every $B_i$ is a principal up-set
$\{y:y_{S_i}=\one\}$ under a product measure. The lower bound is the
Harris--Kleitman inequality (the events $\neg B_i$ are decreasing, hence positively
correlated) \cite{harris1960}. The upper bounds are Janson's inequality
\cite{janson1990} in the product form of the Boppana--Spencer proof
(\cref{thm:janson-general}) and in the form of \cite[Thm.~8.1.1]{alonspencer}.
$D(B_i)$ and $D(B_i\wedge B_j)$ are products of marginals over $S_i$ and
$S_i\cup S_j$.
\end{proof}

So for unate CNFs with $\eps_0\le1/2$ whose dependent pairs carry total probability
$\Delta\le\eps$, the geometric mean of the two ends has relative error at most
$e^{\Delta}-1$, however small $p$ is. Without a bound on $\eps_0$ the last form can
be arbitrarily loose, but the product form, with clauses ordered by decreasing
$b_i$, is at most $M\exp(2\Delta/(1-\Delta))$ for $\Delta<1$: if clause $i$ shares a
variable with an earlier clause $j$, then $d_i\ge D(B_i\wedge B_j)\ge b_i+b_j-1\ge
2b_i-1$, so $1-b_i\ge(1-\Delta)/2$; otherwise $d_i=0$. Example: $m=6\cdot10^4$
random clauses of width $12$ over $2\cdot10^5$ variables give the certified interval
$p\in[4.33,5.07]\cdot10^{-7}$ (ratio $1.17$) in $50$ seconds of pure Python
(wall-clock, on a shared 4-core machine; \cref{tab:janson}). The
version for unate $\wedge$'s of arbitrary monotone sub-circuits, with $b_i$ and
$\Delta$ estimated by sampling, is \cref{thm:janson-general}.
